\section*{Repair sketch: admitted whole-support quotient at nodes [43]--[46]}

\paragraph{Frozen defect.}
The failed implementation treated the node-[43] payload as an arbitrary raw
map on curvature labels and then opened a residual when no closed
representative was supplied separately.  This is not the incoming manuscript
state.  A curvature determination certificate already contains a
\emph{functional admissible rank quotient}
\[
q:\rho^{\rm ex}_{\varnothing}(G)\longrightarrow Q.
\]
By Definition~\texttt{def:admissible-rank-quotient}, admissibility includes
the response-preservation audit and the following conditional certificate:
if the code of (q) is non-injective on the declared finite coordinate
family, it constructs a strictly smaller admissible closed representative.

\paragraph{Directed state.}
The incoming path is the CT15 rank-drop output followed by the proper-support
classification.  Nodes [37], [39], and [42] have already consumed respectively
the target-defective, proper-atom compression, and proper enlarged-support
outcomes.  The only payload at node [43] is therefore the fixed finite response
system, one admitted quotient (q), and two distinct declared coordinates
(a\ne b) identified by its code.  The statement (Z=G) records only the
location of this certificate's finite determination support; it does not
replace the local response system by a universe of ambient graphs or contexts.

\paragraph{Both-sides test.}
There is one finite predicate: injectivity of the admitted code on its declared
coordinate family.
\begin{center}
\begin{tabular}{p{0.22\textwidth}p{0.68\textwidth}}
injective & The purported rank-drop witness (a\ne b) and
  (q(a)=q(b)) are contradictory; this is C5.\\
non-injective & The representative clause already stored in the admitted
  quotient produces a certified strict reduction; minimality closes by C2.\\
\end{tabular}
\end{center}
Both outcomes close.  There is no third ``missing representative'' outcome,
because removing the representative clause changes an admitted quotient back
into a raw proposal and hence changes the type of the predecessor payload.

\paragraph{CT worksheet.}
The CT15 rank-drop witness is routed as closed whole-object data to the exact
whole-support barrier.  S-Def is the declared finite curvature-coordinate
enumeration and its exact response system.  S-Dich is injective versus
non-injective code.  S-Rout transports the same admitted quotient and the two
identified coordinates.  S-Trig is the closed representative clause already
inside admissibility.  The positive leaf is C5 and the negative leaf is C2.
No restoration or recursive measure is needed.

\paragraph{Repair-network identity.}
Node [44] is independent local bookkeeping on one supplied delayed
compensation component.  The graph certificate records a finite connected
component, its boundary leaves, degree-three lower bounds, surplus, and edge
count.  The graph layer derives the handshake and cycle-rank equations and the
core arithmetic theorem derives
\[
s=p-2+2\beta_Z-\sigma_Z.
\]
This calculation neither selects the whole-support quotient nor supplies an
extra hypothesis to its closure.

\paragraph{Practicality and ownership.}
The quotient barrier inspects one proof-carrying admitted certificate and has
constant logical work.  CT15 scans only the explicitly enumerated local wedge
coordinates and has the registered linear work bound.  The repair identity
uses finite degree sums on one supplied component.  No family of graphs,
subgraphs, completions, contexts, quotient proposals, or representatives is
generated.  The admission/injectivity theorem belongs to CT15, the finite
component identity belongs to Graph and Core, and the Erd\H{o}s layer supplies
only the concrete wedge response system and composes the inherited payload.
